\subsection{Equation of a Circle Given a Tangent Line} 
We can find the equation of a circle given a center and a tangent line.
\begin{Example}
Give an equation of the circle with center $(2,-5)$ and with tangent line $y=\frac{2}{3}x+11$.
\begin{itemize}
\item $(h,k)=\blankpair$
\item The radius is the distance from the center to the line.
\end{itemize}
\begin{lpic}{}{9}
\regtextleft{1}{-1}{Perpendicular line:};
\regtextleft{1}{-2}{Find the point of intersection};
\regtextleft{1}{-8}{Point of intersection is};
\regtextleft{1}{-9}{$r^2={}$};
\end{lpic}
\begin{itemize}
\item So the equation of the circle is \makeblank.
\end{itemize}
\end{Example}
This is easier if the line is \makeblank[1in] or \makeblank[1in].
\begin{Example}
Give an equation of the circle with center $(-3,1)$ and tangent to the line $x=2$.
\begin{itemize}
\item $(h,k)=\blankpair$
\item $r$ is the distance from \blankpair to \makeblank[1in]
\item So $r=\makeblank[1in]=\makeblanksmall$
\item The equation of the circle is \makeblank.
\end{itemize}
\end{Example}
The $x$- and $y$-axis are special horizontal and vertical lines.
\begin{Example}
Give an equation of the circle with center $(-1,4)$ and tangent to the $x$-axis.
\begin{itemize}
\item The $x$-axis is the line \makeblank[1in]
\item $(h,k)=\blankpair$
\item $r$ is the distance from \blankpair to \makeblank[1in]
\item So $r=\makeblank[1in]=\makeblanksmall$
\item The equation of the circle is \makeblank.
\end{itemize}
\end{Example}